Le lancer de rayons (\textit{ray tracing}) constitue aujourd'hui une approche numérique largement utilisée dans des domaines scientifiques et industriels variés. Initialement popularisée par la synthèse d'images pour le rendu photoréaliste, cette méthode est désormais employée dans des applications aussi diverses que la simulation optique, le transport de particules ou le transfert radiatif. Malgré la diversité des phénomènes physiques considérés, ces applications reposent sur une problématique géométrique commune : déterminer efficacement les intersections entre un très grand nombre de rayons et une géométrie tridimensionnelle complexe.

Cette problématique a conduit au développement de bibliothèques logicielles spécialisées dans la navigation géométrique et l'accélération des calculs d'intersection. En physique des hautes énergies, le framework \textsc{ROOT} développé au CERN intègre ainsi le module de géométrie \texttt{TGeo}, conçu pour le suivi de trajectoires au sein de détecteurs complexes. Cette infrastructure géométrique a notamment servi de fondation au développement de \textsc{ROBAST}, une bibliothèque de lancer de rayons destinée à la simulation optique des télescopes Cherenkov \cite{Brun1997ROOT,Okumura2016}.

Parallèlement, les progrès réalisés dans le domaine du rendu graphique ont conduit au développement de bibliothèques dédiées à l'accélération des calculs d'intersection rayon--géométrie. Parmi celles-ci, \textsc{Intel Embree} constitue aujourd'hui une référence pour les architectures CPU. Cette bibliothèque fournit des structures d'accélération hiérarchiques (\textit{Bounding Volume Hierarchies}, BVH) ainsi que des algorithmes de parcours fortement optimisés permettant de réduire significativement le coût des calculs géométriques \cite{Wald2014Embree}.

Les performances offertes par ces outils ont conduit à leur adoption dans plusieurs codes scientifiques. Par exemple, le modèle de transfert radiatif tridimensionnel \textsc{DART} s'appuie désormais sur Embree afin d'accélérer les tests d'intersection rayon--triangle lors des simulations de scènes urbaines et forestières complexes \cite{Qi2019HybridScene}. Cette convergence illustre le caractère transversal du lancer de rayons : indépendamment de la physique considérée, la recherche d'intersections constitue une brique algorithmique commune pouvant être mutualisée.

Le code de Monte Carlo Ray Tracing présenté dans ce chapitre s'inscrit dans cette démarche. Il s'appuie sur la bibliothèque Embree pour le traitement géométrique des trajectoires des rayons, tandis que les modèles physiques propres au transfert radiatif sont utilisés pour estimer les facteurs de vue entre les surfaces du domaine.